\chapter{Podsumowanie}

Podsumowując porównanie obu metod XAI, Occlusion okazała się bardziej stabilna przestrzennie, szczególnie w warunkach losowego szumu Gaussa. Jej regionalny sposób wyznaczania istotności powodował, że drobne zmiany wartości pikseli nie prowadziły do gwałtownej zmiany układu map. Było to szczególnie widoczne dla zbioru MNIST, gdzie najważniejsze obszary pozostawały związane z kształtem cyfry nawet przy wysokich poziomach szumu. Dla CIFAR-10 stabilność była niższa, ale Occlusion nadal zachowywała bardziej spójny obraz istotnych regionów niż Integrated Gradients.

Integrated Gradients okazało się bardziej wrażliwe na perturbacje obrazu wejściowego. Wynikało to z gradientowego charakteru tej metody, która reagowała na zmiany lokalnej struktury obrazu oraz na zmiany kierunku gradientów modelu. W przypadku losowego szumu prowadziło to do stopniowego osłabienia zgodności map, natomiast w przypadku FGSM wpływ był silniejszy, ponieważ sama perturbacja była wyznaczana na podstawie gradientu funkcji straty. Szczególnie dla CIFAR-10 mapy Integrated Gradients szybko traciły stabilność, co wskazywało, że szczegółowe wyjaśnienia gradientowe były bardziej podatne na zakłócenia danych wejściowych.

Wyniki nie wskazywały jednak na jednoznaczną przewagę jednej metody we wszystkich aspektach. Occlusion była bardziej odporna pod względem zachowania położenia istotnych regionów, ale Integrated Gradients w części eksperymentów, zwłaszcza dla czystych obrazów CIFAR-10, lepiej wskazywało fragmenty funkcjonalnie istotne dla decyzji modelu. Oznaczało to, że stabilność przestrzenna mapy i skuteczność wskazywania elementów krytycznych dla predykcji nie zawsze były równoważne.

Ostatecznie przeprowadzone eksperymenty pokazały, że stabilność metod XAI nie była cechą stałą, lecz zależała od rodzaju danych, typu perturbacji oraz mechanizmu działania danej metody. Occlusion można było uznać za metodę bardziej odporną przestrzennie na losowe zakłócenia, zwłaszcza dla prostszych danych. Integrated Gradients dostarczało bardziej szczegółowych wyjaśnień, jednak ta szczegółowość wiązała się z większą podatnością na perturbacje, zwłaszcza na perturbacje gradientowe. Najtrudniejszym scenariuszem dla obu metod był atak FGSM na zbiorze CIFAR-10, gdzie złożoność obrazów i ukierunkowany charakter perturbacji prowadziły do wyraźnej destabilizacji map atrybucji.